% ------------------------------------------------------------
% Page 72
% ------------------------------------------------------------

\begin{center}
    {\Large\bfseries LA CARDERIE DE BOURRES DE SOIE}
\end{center}

\vspace{0.55cm}

Ce n’est pas sans raisons que \textbf{Jean-Louis Laget} s’était lancé dans ces travaux,
probablement en deux tranches : la partie Sud, jusqu’à la salle d’entrée, puis l’aile
Nord venant rejoindre le foulon (le Rouilladou), tour couverte d’ardoises..

\vspace{0.35cm}

C’est qu’en effet on était en pleine extension de la culture des vers à soie. Il y avait à
Meyrueis de nombreux mûriers et plusieurs élevages contribuaient à l’aisance de
certaines familles. D’où, probablement, l’idée de monter une petite fabrique de
cardage de bourres de soie : les bourres, parties externes du cocon en soie grossière
ainsi que les déchets du dévidage étaient soigneusement récupérés. La fabrique était
alimentée probablement par des bourres de soie venant également des basses
Cévennes, Le Vigan, Anduze, Saint André de Valborgne... et Meyrueis ! (Le registre
des patentes en donne la description suivante) :

\vspace{0.25cm}

\begin{flushleft}
\hspace{0.35cm}
\parbox{0.88\textwidth}{
\textit{
« Le bâtiment principal est mieux bâti et plus important à l’intérieur qu’il ne
paraît l’être du dehors. Il est divisé en quatre compartiments : les deux
principaux renferment les 7 métiers mécaniques, l’un des autres sert de magasin
pour les fris ons et de logement au contremaître et le quatrième qui n’est pour
ainsi dire qu’un hangar sert pour l’étendage des fris ons de filoselle.
Il y a de plus, attenant à la maison une cour assez vaste qui sert aussi pour
l’étendage.
La marchandise, une fois dégrossie est transportée à Meyrueis dans les deux
vastes salles que l’industriel y tient à sa disposition pour être soumise à l’action
des presses. »
}
}
\end{flushleft}

\vspace{0.25cm}

Le contremaître n’est autre qu’un certain Benjamin Avesque à partir de 1850 et
\textit{« sa maison est assez vaste et assez bien bâtie. Les 7 métiers à carder les bourres de soie sont
mus par une chute d’eau de la force de 6 chevaux vapeur actionnant la grande roue
verticale de 5m de diamètre. Cet établissement ne chôme jamais par manque d’eau. »}
(3 métiers servaient à dégrossir les bourres de soie, 4 autres à les « finir », à quoi il faut
ajouter 20 presses à bras.)

\vspace{0.35cm}

Les bourres de soie, une fois cardées et réduites en fins flocons entraient probablement
dans la composition du feutre pour les chapeaux fabriqués à Meyrueis; ou garnissaient
les édredons. On peut penser que le feutre après son passage au foulon ( l’actuel
Rouilladou ) alimentait en partie la Manufacture de Chapeaux de Feutre qui appartenait
aux Couderc frères de mon arrière grand père Louis Couderc dont la fabrique était
située dans les maisons entre la Placette et le Pont Vieux. Leurs descendants sont les
Verdier de Cambonéral à Saint Jean du Gard. Les bourres de soie filées donnaient la
filoselle avec laquelle on fabriquait des gants, des rideaux ,des bas ou des foulards moins
chers qu’en soie naturelle.
